\section{Expression DSL}\label{sec:expr-dsl}

Agent workflows must evaluate conditions over structured data: tool responses
encoded as JSON values, user-provided parameters, and intermediate computation
results. Naively using Python's \identifiable{eval()} on such data is
impractical because it offers no resource bounds (a malicious or deeply nested
expression can diverge), no totality guarantees (partial functions are hard to
reason about), and no type safety (arbitrary Python code can escape the sandbox).
We address these shortcomings with a purpose-built expression DSL whose
evaluation is governed by big-step operational semantics that guarantee
termination and totality. The evaluator is implemented in OCaml and embedded in
the P-A-R runtime as a total function from expressions and a variable context to
values.

\subsection{Syntax}\label{sec:expr-syntax}

Figure~\ref{fig:expr-grammar} presents the expression grammar side-by-side with
its OCaml implementation. The left side defines 13 expression forms
mathematically; the right side shows the corresponding algebraic data type, which
adds two relational operators (\identifiable{greater\_or\_equal} and
\identifiable{less\_or\_equal}) as implementation extensions.

\begin{figure}[t]
\begin{minipage}{0.48\textwidth}
\small
\[
\begin{array}{rcl}
e &::=& \identifiable{literal}(v) \\
  &\mid& \identifiable{variable}(x) \\
  &\mid& \identifiable{equals}(e_1, e_2) \\
  &\mid& \identifiable{not\_equals}(e_1, e_2) \\
  &\mid& \identifiable{greater\_than}(e_1, e_2) \\
  &\mid& \identifiable{less\_than}(e_1, e_2) \\
  &\mid& \identifiable{and}(e_1, e_2) \\
  &\mid& \identifiable{or}(e_1, e_2) \\
  &\mid& \identifiable{not}(e) \\
  &\mid& \identifiable{contains}(e_1, e_2) \\
  &\mid& \identifiable{has\_key}(e, k) \\
  &\mid& \identifiable{is\_empty}(e) \\
  &\mid& \identifiable{matches}(e, p)
\end{array}
\]
\end{minipage}
\hfill
\begin{minipage}{0.48\textwidth}
\small
\begin{lstlisting}[style=ocaml]
type expr =
  | Literal of json
  | Variable of string
  | Equals of expr * expr
  | Not_equals of expr * expr
  | Greater_than of expr * expr
  | Less_than of expr * expr
  | Greater_or_equal of expr * expr
  | Less_or_equal of expr * expr
  | And of expr * expr
  | Or of expr * expr
  | Not of expr
  | Contains of expr * expr
  | Has_key of expr * string
  | Is_empty of expr
  | Matches of expr * string
\end{lstlisting}
\end{minipage}
\caption{Expression grammar (left) and OCaml implementation (right).}
\label{fig:expr-grammar}
\end{figure}

We write $\Gamma \vdash e \downarrow v$ to mean that expression $e$ evaluates to
value $v$ under variable context $\Gamma$, which maps variable names (dot-notation
paths, e.g.\ \texttt{data.items[0].name}) to JSON values.

\subsection{Big-Step Semantics}\label{sec:expr-semantics}

The evaluation relation is defined by the inference rules below, organised by
operational category.

\medskip\noindent
\textbf{Literal and variable:}

\begin{mathpar}
\inferrule{ }{\Gamma \vdash \identifiable{literal}(v) \downarrow v}
\and
\inferrule{\Gamma(x) = v}{\Gamma \vdash \identifiable{variable}(x) \downarrow v}
\end{mathpar}

\medskip\noindent
\textbf{Equality predicates:}

\begin{mathpar}
\inferrule{
  \Gamma \vdash e_1 \downarrow v_1 \quad
  \Gamma \vdash e_2 \downarrow v_2 \quad
  v_1 =_{\identifiable{JSON}} v_2
}{\Gamma \vdash \identifiable{equals}(e_1, e_2) \downarrow \identifiable{true}}
\and
\inferrule{
  \Gamma \vdash e_1 \downarrow v_1 \quad
  \Gamma \vdash e_2 \downarrow v_2 \quad
  v_1 \neq_{\identifiable{JSON}} v_2
}{\Gamma \vdash \identifiable{not\_equals}(e_1, e_2) \downarrow \identifiable{true}}
\end{mathpar}

\medskip\noindent
\textbf{Relational predicates:}

\begin{mathpar}
\inferrule{
  \Gamma \vdash e_1 \downarrow v_1 \quad
  \Gamma \vdash e_2 \downarrow v_2 \quad
  v_1 >_{\identifiable{JSON}} v_2
}{\Gamma \vdash \identifiable{greater\_than}(e_1, e_2) \downarrow \identifiable{true}}
\and
\inferrule{
  \Gamma \vdash e_1 \downarrow v_1 \quad
  \Gamma \vdash e_2 \downarrow v_2 \quad
  v_1 <_{\identifiable{JSON}} v_2
}{\Gamma \vdash \identifiable{less\_than}(e_1, e_2) \downarrow \identifiable{true}}
\and
\inferrule{
  \Gamma \vdash e_1 \downarrow v_1 \quad
  \Gamma \vdash e_2 \downarrow v_2 \quad
  v_1 \geq_{\identifiable{JSON}} v_2
}{\Gamma \vdash \identifiable{greater\_or\_equal}(e_1, e_2) \downarrow \identifiable{true}}
\and
\inferrule{
  \Gamma \vdash e_1 \downarrow v_1 \quad
  \Gamma \vdash e_2 \downarrow v_2 \quad
  v_1 \leq_{\identifiable{JSON}} v_2
}{\Gamma \vdash \identifiable{less\_or\_equal}(e_1, e_2) \downarrow \identifiable{true}}
\end{mathpar}

\medskip\noindent
\textbf{Logical connectives:}

\begin{mathpar}
\inferrule{
  \Gamma \vdash e_1 \downarrow \identifiable{true} \quad
  \Gamma \vdash e_2 \downarrow \identifiable{true}
}{\Gamma \vdash \identifiable{and}(e_1, e_2) \downarrow \identifiable{true}}
\and
\inferrule{
  \Gamma \vdash e_1 \downarrow \identifiable{false}
}{\Gamma \vdash \identifiable{and}(e_1, e_2) \downarrow \identifiable{false}}
\and
\inferrule{
  \Gamma \vdash e_1 \downarrow \identifiable{true}
}{\Gamma \vdash \identifiable{or}(e_1, e_2) \downarrow \identifiable{true}}
\and
\inferrule{
  \Gamma \vdash e_1 \downarrow \identifiable{false} \quad
  \Gamma \vdash e_2 \downarrow v_2
}{\Gamma \vdash \identifiable{or}(e_1, e_2) \downarrow v_2}
\and
\inferrule{
  \Gamma \vdash e \downarrow \identifiable{false}
}{\Gamma \vdash \identifiable{not}(e) \downarrow \identifiable{true}}
\and
\inferrule{
  \Gamma \vdash e \downarrow \identifiable{true}
}{\Gamma \vdash \identifiable{not}(e) \downarrow \identifiable{false}}
\end{mathpar}

\medskip\noindent
\textbf{Collection operations:}

\begin{mathpar}
\inferrule{
  \Gamma \vdash e_1 \downarrow v_1 \quad
  \Gamma \vdash e_2 \downarrow v_2 \quad
  \identifiable{contains}(v_1, v_2) = \identifiable{true}
}{\Gamma \vdash \identifiable{contains}(e_1, e_2) \downarrow \identifiable{true}}
\and
\inferrule{
  \Gamma \vdash e \downarrow v \quad
  \identifiable{has\_key}(v, k) = \identifiable{true}
}{\Gamma \vdash \identifiable{has\_key}(e, k) \downarrow \identifiable{true}}
\and
\inferrule{
  \Gamma \vdash e \downarrow v \quad
  \identifiable{is\_empty}(v) = \identifiable{true}
}{\Gamma \vdash \identifiable{is\_empty}(e) \downarrow \identifiable{true}}
\end{mathpar}

\medskip\noindent
\textbf{Regex matching:}

\begin{mathpar}
\inferrule{
  \Gamma \vdash e \downarrow v \quad
  \identifiable{matches}(v, p) = \identifiable{true}
}{\Gamma \vdash \identifiable{matches}(e, p) \downarrow \identifiable{true}}
\end{mathpar}

All remaining cases (where the predicate does not hold) evaluate to
$\identifiable{false}$ by symmetry with the rules above.

\subsection{Resource Bounds and Properties}\label{sec:expr-properties}

The evaluator enforces two resource limits that make termination formal:

\begin{prop}[Resource Limits]
\label{prop:resource-limits}
The evaluator maintains a visit counter bounded by
\identifiable{max\_node\_visits} = 1000 and a recursion depth bounded by
\identifiable{max\_depth} = 10.
\end{prop}

\begin{prop}[Termination]
\label{prop:termination}
For any expression $e$ and context $\Gamma$, the evaluator terminates in at most
1000 node visits (validated empirically in \S\ref{sec:evaluation}).
\end{prop}

\begin{prop}[Totality]
\label{prop:totality}
The evaluator is a total function: for all $e$ and $\Gamma$, either
$\Gamma \vdash e \downarrow v$ for some value $v$, or the evaluator raises
\identifiable{Resource\_limit} which is caught and converted to an error result
(validated empirically in \S\ref{sec:evaluation}).
\end{prop}

The totality property follows from structural recursion over the expression tree
with a decremented visit counter; when the counter reaches zero the evaluator
raises \identifiable{Resource\_limit} rather than diverging.

\subsection{Example Trace}\label{sec:expr-example}

Consider evaluating the expression%
\[
\identifiable{and}(
  \identifiable{equals}(
    \identifiable{variable}(\text{\texttt{"status"}}),
    \identifiable{literal}(\text{\texttt{"running"}})),
  \identifiable{not}(
    \identifiable{is\_empty}(
      \identifiable{variable}(\text{\texttt{"results"}}))))
\]
under the context $\Gamma = \{\,\text{\texttt{"status"}} \mapsto
\text{\texttt{"running"}},\ \text{\texttt{"results"}} \mapsto [1,2,3]\,\}$.

\begin{enumerate}
\item $\Gamma \vdash \identifiable{variable}(\text{\texttt{"status"}}) \downarrow
  \text{\texttt{"running"}}$ \quad (lookup in $\Gamma$)
\item $\Gamma \vdash \identifiable{literal}(\text{\texttt{"running"}}) \downarrow
  \text{\texttt{"running"}}$ \quad (literal rule)
\item $\Gamma \vdash \identifiable{equals}(\cdots) \downarrow \identifiable{true}$
  \quad (equality holds: both sides are \texttt{"running"})
\item $\Gamma \vdash \identifiable{variable}(\text{\texttt{"results"}}) \downarrow
  [1,2,3]$ \quad (lookup in $\Gamma$)
\item $\Gamma \vdash \identifiable{is\_empty}([1,2,3]) \downarrow \identifiable{false}$
  \quad (collection is non-empty)
\item $\Gamma \vdash \identifiable{not}(\identifiable{false}) \downarrow
  \identifiable{true}$ \quad (negation rule)
\item $\Gamma \vdash \identifiable{and}(\identifiable{true},
  \identifiable{true}) \downarrow \identifiable{true}$ \quad (both conjuncts true)
\end{enumerate}
The short-circuit semantics of \identifiable{and} means step 5 is reached only
after step 3 succeeds; similarly, step 6 is reached only after step 4 yields
$\identifiable{false}$, at which point the negation flips it to $\identifiable{true}$.
The trace uses 7 node visits, well within the 1000-visit bound.